% year: תשפ"ג
% type: מבחן
% semester: א
% moed: א
% number: 6א
% points: 10
% categories: integral-applications
% source: exams/מבחנים/תשפג/מועד א תשפג.pdf

\begin{question}
חשבו את השטח החסום בין הגרפים של הפונקציות $f(x)=\sin^2x$ ו-$g(x)=\frac14$ בקטע $[0,\pi]$.
\end{question}

\begin{hint}
מצאו את נקודות החיתוך $\sin^2x=\frac14$ בקטע, וחלקו את הקטע לפי איזו פונקציה גדולה יותר. השתמשו בזהות $\sin^2x=\frac{1-\cos 2x}{2}$.
\end{hint}

\begin{solution}
\textbf{נקודות חיתוך:} $\sin^2x=\frac14\iff\sin x=\pm\frac12$. בקטע $[0,\pi]$ מתקיים $\sin x\ge 0$, ולכן $\sin x=\frac12$, כלומר $x=\frac{\pi}{6}$ או $x=\frac{5\pi}{6}$.

בקטעים $[0,\frac{\pi}{6})$ ו-$(\frac{5\pi}{6},\pi]$ מתקיים $0\le\sin x<\frac12$, ולכן $\sin^2x<\frac14$; בקטע $(\frac{\pi}{6},\frac{5\pi}{6})$ מתקיים $\sin^2x>\frac14$.

\textbf{פונקציה קדומה:} לפי $\sin^2x=\frac{1-\cos2x}{2}$,
\[ \int\left(\sin^2x-\frac14\right)dx=\int\left(\frac14-\frac{\cos 2x}{2}\right)dx=\frac{x}{4}-\frac{\sin 2x}{4}=:F(x). \]
ערכים: $F(0)=0$, $F(\frac{\pi}{6})=\frac{\pi}{24}-\frac{\sqrt3}{8}$, $F(\frac{5\pi}{6})=\frac{5\pi}{24}+\frac{\sqrt3}{8}$, $F(\pi)=\frac{\pi}{4}$.

\textbf{השטח} הוא $\int_0^\pi|f-g|\,dx$:
\begin{align*}
S_1&=\int_0^{\pi/6}\left(\tfrac14-\sin^2x\right)dx=-\big(F(\tfrac{\pi}{6})-F(0)\big)=\frac{\sqrt3}{8}-\frac{\pi}{24},\\
S_2&=\int_{\pi/6}^{5\pi/6}\left(\sin^2x-\tfrac14\right)dx=F(\tfrac{5\pi}{6})-F(\tfrac{\pi}{6})=\frac{\pi}{6}+\frac{\sqrt3}{4},\\
S_3&=\int_{5\pi/6}^{\pi}\left(\tfrac14-\sin^2x\right)dx=-\big(F(\pi)-F(\tfrac{5\pi}{6})\big)=\frac{\sqrt3}{8}-\frac{\pi}{24}.
\end{align*}
סה"כ:
\[ S=S_1+S_2+S_3=\frac{\sqrt3}{2}+\frac{\pi}{12}\approx 1.128. \]
(אם מתכוונים רק לתחום הסגור שבין שתי נקודות החיתוך, שטחו $S_2=\frac{\pi}{6}+\frac{\sqrt3}{4}\approx 0.957$.)
\end{solution}
